\subsection{Метод приведённого градиента Вульфа}
\label{subsec:reduced-gradient}

\[
  \varphi(x) \to \min
  \qquad\text{при}\qquad
  A x = b , \quad x \ge 0 ,
\]
где $A$ — матрица $m \times n$, $\rank A = m$, $\varphi$ дифференцируема.
Метод строит последовательность допустимых точек $x^0, x^1, \dots$ с
$Ax^k = b$, $x^k \ge 0$ и убывающим значением цели.

\subsubsection*{Разбиение на базисные и небазисные переменные}

Выберем $m$ линейно независимых столбцов $A$; они образуют базисную матрицу
$B$, остальные — матрицу $N$. Соответственно
\[
  x = \begin{pmatrix} x_B \\ x_N \end{pmatrix} ,
  \qquad
  Ax = b \iff B x_B + N x_N = b
  \qquad\Longrightarrow\qquad
  x_B = B^{-1}\big(b - N x_N\big) .
\]
Базисные переменные выражаются через небазисные — это и есть «приведение»
задачи.

\begin{lemma}\label{lem:wolfe-feasible}
Если $Ax = b$ и $Ad = 0$, то $A(x + \alpha d) = b$ при любом $\alpha$.
\end{lemma}

\begin{proof}
$A(x + \alpha d) = Ax + \alpha A d = b + \alpha\cdot 0 = b$.
\end{proof}

Разбивая направление так же, $d = (d_B;\ d_N)$, из $Ad = 0$ получаем
$B d_B + N d_N = 0$, то есть
\begin{equation}\label{eq:wolfe-dB}
  d_B = -B^{-1} N d_N .
\end{equation}
Значит достаточно выбрать небазисную часть $d_N$, базисная определяется
однозначно.

\subsubsection*{Приведённый градиент}

Пусть $g = \nabla\varphi(x) = (g_B;\ g_N)$. Вектор
\[
  r_N = g_N - N\T B^{-\mathsf{T}} g_B
\]
называется \textbf{приведённым градиентом}.

\begin{theorem}\label{thm:reduced-grad}
Для всякого направления $d$ с $Ad = 0$ выполняется
$g\T d = r_N\T d_N$.
\end{theorem}

\begin{proof}
По~\eqref{eq:wolfe-dB}
\[
  g\T d = g_B\T d_B + g_N\T d_N
        = g_B\T\big(-B^{-1}N d_N\big) + g_N\T d_N
        = \big(g_N\T - g_B\T B^{-1}N\big) d_N .
\]
Транспонируя скобку, $g_N\T - g_B\T B^{-1}N
= \big(g_N - N\T B^{-\mathsf{T}} g_B\big)\T = r_N\T$, откуда
$g\T d = r_N\T d_N$.
\end{proof}

\begin{idea}
Смысл: после исключения базисных переменных задача зависит только от $x_N$, и
$r_N$ — градиент этой «приведённой» функции. Первое слагаемое $g_N$ — прямое
влияние небазисных переменных, второе $N\T B^{-\mathsf T}g_B$ — косвенное,
через вынужденный пересчёт $x_B$. Если ограничений-равенств нет ($m = 0$), то
$r_N = g$, и метод вырождается в обычный градиентный спуск с учётом $x \ge 0$.
\end{idea}

\subsubsection*{Направление Вульфа}

Если небазисная переменная не на границе, можно двигаться против приведённого
градиента; если она на границе $x_j = 0$, отрицательная компонента $d_j$
недопустима. Отсюда правило
\begin{equation}\label{eq:wolfe-dir}
  d_j =
  \begin{cases}
    -r_j , & x_j > 0 , \\
    -r_j , & x_j = 0,\ r_j < 0 , \\
    0 , & x_j = 0,\ r_j \ge 0 ,
  \end{cases}
  \qquad j \ \text{— небазисный},
\end{equation}
а базисная часть берётся из~\eqref{eq:wolfe-dB}.

\begin{theorem}\label{thm:wolfe-descent}
Если направление построено по правилу~\eqref{eq:wolfe-dir} и $d_N \ne 0$, то
$g\T d < 0$, то есть $d$ — направление спуска.
\end{theorem}

\begin{proof}
По теореме~\ref{thm:reduced-grad} $g\T d = r_N\T d_N = \sum_j r_j d_j$.
Разберём слагаемые:
\begin{itemize}[nosep,left=0pt]
  \item $x_j > 0$: $d_j = -r_j$, значит $r_j d_j = -r_j^2 \le 0$;
  \item $x_j = 0$, $r_j < 0$: $d_j = -r_j > 0$, снова $r_j d_j = -r_j^2 < 0$;
  \item $x_j = 0$, $r_j \ge 0$: $d_j = 0$, слагаемое нулевое.
\end{itemize}
Все слагаемые неположительны, а при $d_N \ne 0$ хотя бы одно строго
отрицательно. Значит $g\T d < 0$.
\end{proof}

\subsubsection*{Максимальный допустимый шаг}

Новая точка $x(\alpha) = x + \alpha d$; равенства сохраняются по
лемме~\ref{lem:wolfe-feasible}, остаётся неотрицательность. Для компонент с
$d_j < 0$ нужно $x_j + \alpha d_j \ge 0$, то есть $\alpha \le x_j/(-d_j)$,
откуда
\[
  \alpha_{\max} = \min_{j :\, d_j < 0} \frac{x_j}{-d_j} ,
\]
а если $d_j \ge 0$ при всех $j$, то $\alpha_{\max} = +\infty$.

\begin{theorem}\label{thm:wolfe-step}
Если $x \ge 0$, $Ad = 0$ и $0 \le \alpha \le \alpha_{\max}$, то
$x + \alpha d \ge 0$ и $A(x + \alpha d) = b$.
\end{theorem}

\begin{proof}
Равенства — лемма~\ref{lem:wolfe-feasible}. Для $d_j \ge 0$ очевидно
$x_j + \alpha d_j \ge x_j \ge 0$. Для $d_j < 0$ по определению $\alpha_{\max}$
имеем $\alpha \le x_j/(-d_j)$; умножая на $-d_j > 0$, получаем
$-\alpha d_j \le x_j$, то есть $x_j + \alpha d_j \ge 0$.
\end{proof}

\subsubsection*{Условие оптимальности}

\begin{theorem}\label{thm:wolfe-opt}
Пусть $\varphi$ выпукла и дифференцируема, $x$ допустима и для всех
небазисных индексов
\[
  x_j > 0 \implies r_j = 0 ,
  \qquad
  x_j = 0 \implies r_j \ge 0 .
\]
Тогда $x$ — глобальный минимум задачи.
\end{theorem}

\begin{proof}
Возьмём любую допустимую $y$. Так как $Ax = b$ и $Ay = b$, вектор
$d = y - x$ удовлетворяет $Ad = 0$, и по теореме~\ref{thm:reduced-grad}
$g\T(y - x) = r_N\T d_N$. Разберём слагаемые: при $x_j > 0$ имеем $r_j = 0$ и
вклад нулевой; при $x_j = 0$ из допустимости $y_j \ge 0$ следует
$d_j = y_j - x_j = y_j \ge 0$, а $r_j \ge 0$, значит $r_j d_j \ge 0$.
Следовательно $g\T(y-x) \ge 0$. По неравенству первого порядка для выпуклой
функции
\[
  \varphi(y) \ \ge\ \varphi(x) + g\T(y - x) \ \ge\ \varphi(x) .
\]
Это верно для любой допустимой $y$, значит $x$ — глобальный минимум.
\end{proof}

\begin{idea}
Записанное условие — это в точности условия ККТ задачи
$\varphi(x) \to \min$, $Ax = b$, $x \ge 0$: приведённый градиент играет роль
оценок $d[j]$ симплекс-метода (\S\ref{subsec:simplex-scheme}), а условие
«$x_j > 0 \Rightarrow r_j = 0$» — дополняющей нежёсткости. При линейной
$\varphi$ метод Вульфа превращается в симплекс-метод.
\end{idea}

\begin{algo}[title={Метод приведённого градиента Вульфа}]
\begin{steps}
  \item Задать допустимую $x^0 \ge 0$, $Ax^0 = b$; выбрать базис $B$; задать
        $\eps > 0$; $k = 0$.
  \item Разбить переменные $x^k = (x_B^k;\ x_N^k)$ и вычислить
        $g^k = \nabla\varphi(x^k) = (g_B^k;\ g_N^k)$.
  \item Вычислить приведённый градиент
        $r_N^k = g_N^k - N\T B^{-\mathsf T} g_B^k$.
  \item Построить $d_N^k$ по правилу~\eqref{eq:wolfe-dir}.
  \item Если $d_N^k = 0$ — допустимого направления спуска нет; в выпуклом
        случае $x^k$ оптимальна, конец.
  \item Найти $d_B^k = -B^{-1}N d_N^k$ и полное направление $d^k$.
  \item Вычислить $\alpha_{\max} = \min_{j:\,d_j^k<0} x_j^k / (-d_j^k)$.
  \item Найти $\alpha_k = \argmin_{0 \le \alpha \le \alpha_{\max}}
        \varphi(x^k + \alpha d^k)$.
  \item $x^{k+1} = x^k + \alpha_k d^k$. Если какая-то переменная обратилась в
        нуль, при необходимости сменить базис.
  \item $k := k+1$, перейти к шагу~2.
\end{steps}
\end{algo}

\begin{example}
$\varphi(x) = (x_1-0{,}2)^2 + (x_2-0{,}3)^2 + (x_3-0{,}5)^2 \to \min$ при
$x_1+x_2+x_3 = 1$, $x \ge 0$. Здесь $A = (1\ \ 1\ \ 1)$, $b = 1$. Возьмём
$x^0 = (0{,}6;\,0{,}1;\,0{,}3)\T$ и базисную переменную $x_3$, тогда $B = (1)$,
$N = (1\ \ 1)$ и $d_3 = -(d_1 + d_2)$.

Градиент $g^0 = (0{,}8;\,-0{,}4;\,-0{,}4)\T$, приведённый градиент
$r_N^0 = (0{,}8;\,-0{,}4)\T - (1;\,1)\T(-0{,}4) = (1{,}2;\,0)\T$. Обе небазисные
переменные положительны, значит $d_N^0 = -r_N^0 = (-1{,}2;\,0)\T$,
$d_3^0 = 1{,}2$. Проверка: $Ad^0 = 0$, $(g^0)\T d^0 = -1{,}44 < 0$.
Ограничение даёт $\alpha_{\max} = 0{,}6/1{,}2 = 0{,}5$; одномерная минимизация
$\varphi = 0{,}24 - 1{,}44\alpha + 2{,}88\alpha^2$ даёт $\alpha_0 = 0{,}25$,
откуда $x^1 = (0{,}3;\,0{,}1;\,0{,}6)\T$ и $\varphi$ падает с $0{,}24$ до
$0{,}06$.

Вторая итерация: $g^1 = (0{,}2;\,-0{,}4;\,0{,}2)\T$,
$r_N^1 = (0;\,-0{,}6)\T$, $d^1 = (0;\,0{,}6;\,-0{,}6)\T$,
$\alpha_{\max} = 1$, $\alpha_1 = 0{,}25$, откуда
$x^2 = (0{,}3;\,0{,}25;\,0{,}45)\T$ и $\varphi = 0{,}015$. Итого
$0{,}24 \to 0{,}06 \to 0{,}015$.
\end{example}
